% =====================================================================
%  第二部分  数值方法原理
% =====================================================================
\part{数值方法原理}

第一部分得到了两个偏微分方程和四个条件。这一部分回答一个问题：
这些方程为什么算不出来，以及我们用什么办法算出来。

% =====================================================================
\section{为什么必须用数值方法}

在动手之前，先说明为什么不能像解一元二次方程那样写一个公式就完事。
有四个障碍，一个比一个麻烦。

\subsection{障碍一：扩散系数随含水率变化}

问题一的扩散系数是

\begin{equation}
D = 7\times10^{-9}e^{-0.89/C}.
\end{equation}

这个式子里的 $C$ 是未知的含水率。也就是说，
要求解方程必须先知道 $D$，而 $D$ 又依赖于要求解的 $C$。
这种方程叫非线性方程，不能像线性方程那样把解写成级数叠加的形式。

先感受一下这个公式的行为。当 $C$ 从 $2.55$ 降到 $0.15$，

\begin{equation}
D(2.55) = 7\times10^{-9}e^{-0.89/2.55}=4.94\times10^{-9},
\end{equation}

\begin{equation}
D(0.15) = 7\times10^{-9}e^{-0.89/0.15}=7\times10^{-9}\times e^{-5.933}
= 7\times10^{-9}\times2.65\times10^{-3}=1.86\times10^{-11}.
\end{equation}

两者相差 $266$ 倍。也就是说，越干的地方水分越难移动，
这正好对应干燥后期速度越来越慢的现象。
但这也意味着 $D$ 在空间上是处处不同的，不能提到积分号或导数号外面。

\subsection{障碍二：物性同时依赖两个未知量}

在问题二到问题四中，扩散系数写成

\begin{equation}
D = 2.4\times10^{-3}e^{-0.45/C}e^{-3850/T},
\end{equation}

它同时依赖含水率 $C$ 和温度 $T$。
而温度本身也要通过传热方程求解，传热方程里的导热系数 $k$ 又依赖 $C$。
两个方程互相纠缠，必须联立求解，不能分开各算各的。

\subsection{障碍三：边界条件是耦合的}

表面条件里出现了表面值 $T(R_0,t)$ 和 $C(R_0,t)$，
它们本身是未知的，而且和边界的梯度联系在一起。
这种边界条件不能直接代入，必须和内部的解一起确定。

\subsection{障碍四：求解区域会变形}

问题四里药材半径从 $2.000$ 厘米收缩到 $1.198$ 厘米。
这意味着我们要求解的区域本身在缩小，边界在移动。
这类问题叫移动边界问题，处理起来需要额外的技巧，第八章详细讲。

\keypoint
四个障碍导致这道题没有现成的公式解。
唯一的办法是把连续的区域切成很多小格，
在每个小格上用代数运算近似微分运算，
把微分方程变成一大组代数方程，再交给计算机求解。
这就是数值方法。

% =====================================================================
\section{坐标变换：把麻烦的 1/r 去掉}

这一章讲本套方法最核心的一个技巧。它解决的问题是：
第一章里那个 $1/r$ 在圆心附近会带来麻烦。

\subsection{麻烦在哪里}

回看式 \eqref{eq:lap-expand}，算子可以写成

\begin{equation}
\frac{1}{r}\frac{\partial}{\partial r}\left(r\,D\,\frac{\partial C}{\partial r}\right)
= D\frac{\partial^2 C}{\partial r^2}
+ \frac{D}{r}\frac{\partial C}{\partial r}
+ \frac{\partial D}{\partial r}\frac{\partial C}{\partial r}.
\label{eq:singular}
\end{equation}

中间那一项 $\dfrac{D}{r}\dfrac{\partial C}{\partial r}$ 里有 $1/r$。

当 $r$ 趋向于零时，$\dfrac{\partial C}{\partial r}$ 也趋向于零，
两个无穷小相除，极限是有限的。数学上没有问题。
但是对计算机来说有问题：
计算机处理的是离散的点，在 $r=0$ 这个点上，
它只会看到分母是零，无法直接算。

常规的做法是对圆心附近做特殊处理。
比如把最里面的那个控制体取成一个半径 $h/2$ 的小圆柱，
把它的体积单独算出来，再单独写它对应的方程。
这样做当然可以，但会让程序变得复杂，
而且要额外推导公式，容易出错。

\subsection{换元：用 u 代替 r}

我们采取另一条路。定义一个新的自变量

\begin{equation}
u = \left(\frac{r}{R}\right)^2 = \xi^2,
\label{eq:utrans}
\end{equation}

其中 $\xi=r/R$ 是归一化的半径，取值范围从 $0$ 到 $1$。
新变量 $u$ 的取值范围也是从 $0$ 到 $1$。

当 $r=0$ 时 $u=0$，当 $r=R$ 时 $u=1$。

\why{为什么想到平方}
因为圆柱和球的算子里的 $r$ 总是以 $r$ 乘以梯度的形式出现，
也就是 $r\dfrac{\partial}{\partial r}$。
如果令 $u$ 正比于 $r^2$，那么 $\mathrm{d}u$ 正比于 $r\,\mathrm{d}r$，
这个 $r$ 恰好可以和算子里的 $r$ 配对消掉。
这是一个基于结构的观察，不是碰巧试出来的。

\subsection{第一步：求出两个导数之间的换算关系}

从式 \eqref{eq:utrans} 出发，先反解出 $r$：

\begin{equation}
u = \frac{r^2}{R^2}
\quad\Longrightarrow\quad
r^2 = R^2 u
\quad\Longrightarrow\quad
r = R\sqrt{u}.
\label{eq:rsqrtu}
\end{equation}

对式 \eqref{eq:rsqrtu} 求 $u$ 的导数。回忆 $\dfrac{\mathrm{d}}{\mathrm{d}u}\sqrt{u}
= \dfrac{1}{2\sqrt{u}}$，所以

\begin{equation}
\frac{\partial r}{\partial u} = R\cdot\frac{1}{2\sqrt{u}} = \frac{R}{2\sqrt{u}}.
\label{eq:drdu}
\end{equation}

链式法则告诉我们，如果 $f$ 是 $u$ 的函数，
那么 $\dfrac{\partial f}{\partial u}=\dfrac{\partial f}{\partial r}\cdot\dfrac{\partial r}{\partial u}$，
于是

\begin{equation}
\frac{\partial}{\partial u} = \frac{R}{2\sqrt{u}}\frac{\partial}{\partial r}.
\end{equation}

把上式两边同乘 $\dfrac{2\sqrt{u}}{R}$，得到我们真正要用的形式

\begin{equation}
\boxed{\;\frac{\partial}{\partial r} = \frac{2\sqrt{u}}{R}\frac{\partial}{\partial u}.\;}
\label{eq:drdx}
\end{equation}

\subsection{第二步：逐项改造算子}

现在把式 \eqref{eq:singular} 中的每一项拆开处理。
为书写方便，记算子为

\begin{equation}
\mathcal{L}\phi = \frac{1}{r}\frac{\partial}{\partial r}\left(r\,D\,\frac{\partial\phi}{\partial r}\right),
\end{equation}

其中 $\phi$ 代表 $C$ 或 $T$。

\textbf{第一个环节}，计算 $\dfrac{\partial\phi}{\partial r}$。直接用式 \eqref{eq:drdx}：

\begin{equation}
\frac{\partial\phi}{\partial r} = \frac{2\sqrt{u}}{R}\frac{\partial\phi}{\partial u}.
\label{eq:step1}
\end{equation}

\textbf{第二个环节}，计算 $D\dfrac{\partial\phi}{\partial r}$。把式 \eqref{eq:step1} 乘上 $D$：

\begin{equation}
D\frac{\partial\phi}{\partial r} = D\cdot\frac{2\sqrt{u}}{R}\frac{\partial\phi}{\partial u}.
\label{eq:step2}
\end{equation}

\textbf{第三个环节}，计算 $r\,D\dfrac{\partial\phi}{\partial r}$。
把式 \eqref{eq:rsqrtu} 的 $r=R\sqrt{u}$ 和式 \eqref{eq:step2} 相乘：

\begin{equation}
r\,D\frac{\partial\phi}{\partial r}
= R\sqrt{u}\cdot D\cdot\frac{2\sqrt{u}}{R}\frac{\partial\phi}{\partial u}.
\end{equation}

注意这里出现了 $\sqrt{u}\cdot\sqrt{u}=u$，而且 $R$ 和 $R$ 相消，
于是

\begin{equation}
r\,D\frac{\partial\phi}{\partial r}
= 2\,u\,D\frac{\partial\phi}{\partial u}.
\label{eq:step3}
\end{equation}

这一步是整段推导的关键：原来的一堆根号全部合并成了一个干净的 $u$。

\textbf{第四个环节}，计算 $\dfrac{\partial}{\partial r}\left(rD\dfrac{\partial\phi}{\partial r}\right)$。
把式 \eqref{eq:drdx} 作用到式 \eqref{eq:step3} 上：

\begin{equation}
\frac{\partial}{\partial r}\left(rD\frac{\partial\phi}{\partial r}\right)
= \frac{2\sqrt{u}}{R}\frac{\partial}{\partial u}\left(2uD\frac{\partial\phi}{\partial u}\right).
\end{equation}

\textbf{第五个环节}，除以 $r$。用 $r=R\sqrt{u}$：

\begin{equation}
\frac{1}{r}\frac{\partial}{\partial r}\left(rD\frac{\partial\phi}{\partial r}\right)
= \frac{1}{R\sqrt{u}}\cdot\frac{2\sqrt{u}}{R}\cdot
\frac{\partial}{\partial u}\left(2uD\frac{\partial\phi}{\partial u}\right).
\end{equation}

此时 $\dfrac{1}{\sqrt{u}}\cdot\sqrt{u}=1$，根号再次消掉，剩下

\begin{equation}
\frac{1}{r}\frac{\partial}{\partial r}\left(rD\frac{\partial\phi}{\partial r}\right)
= \frac{2}{R^2}\frac{\partial}{\partial u}\left(2\,u\,D\frac{\partial\phi}{\partial u}\right).
\end{equation}

把括号外的 $2$ 和括号内的 $2$ 合起来，得到

\begin{equation}
\boxed{\;
\frac{1}{r}\frac{\partial}{\partial r}\left(rD\frac{\partial\phi}{\partial r}\right)
= \frac{4}{R^2}\frac{\partial}{\partial u}\left(u\,D\,\frac{\partial\phi}{\partial u}\right).
\;}
\label{eq:reg}
\end{equation}

\subsection{第三步：检查圆心处的表现}

现在看式 \eqref{eq:reg} 的右边在 $u=0$ 处的表现。
把括号里的导数用乘积法则展开：

\begin{equation}
\frac{\partial}{\partial u}\left(uD\frac{\partial\phi}{\partial u}\right)
= D\frac{\partial\phi}{\partial u} + u\frac{\partial}{\partial u}\left(D\frac{\partial\phi}{\partial u}\right).
\label{eq:expand-reg}
\end{equation}

当 $u=0$ 时，上式第一项是 $D\dfrac{\partial\phi}{\partial u}$，
因为 $D$ 在圆心处是一个有限的数，$\dfrac{\partial\phi}{\partial u}$ 也是有限的，
所以第一项有限。第二项里有因子 $u$，在 $u=0$ 处等于零。

两句话合起来：右边在 $u=0$ 处是一个有限的数，没有任何除以零的风险。
麻烦被彻底消掉了。

\why{这为什么比原来的写法好}
原来在 $r$ 坐标下，圆心处的方程需要单独推导，需要知道该处控制体的特殊体积。
现在在 $u$ 坐标下，圆心的方程就是同一个公式，直接把 $u=0$ 代进去就行。
程序的写法因此统一，也不容易出错。

\subsection{第四步：把方程改写成新坐标下的样子}

传质方程原来是

\begin{equation}
\frac{\partial C}{\partial t}=\frac{1}{r}\frac{\partial}{\partial r}\left(rD\frac{\partial C}{\partial r}\right).
\end{equation}

把右边换成式 \eqref{eq:reg} 的形式，得到

\begin{equation}
\boxed{\;
\frac{\partial C}{\partial t}=\frac{4}{R^2}\frac{\partial}{\partial u}\left(uD\frac{\partial C}{\partial u}\right).
\;}
\label{eq:mass-u}
\end{equation}

传热方程同理，

\begin{equation}
\boxed{\;
\rho c_p\frac{\partial T}{\partial t}=\frac{4}{R^2}\frac{\partial}{\partial u}\left(uk\frac{\partial T}{\partial u}\right).
\;}
\label{eq:energy-u}
\end{equation}

\why{为什么右边的未知量还是写成 C 和 T 而不是改写}
$C$ 和 $T$ 是物理量，换坐标不改变它们的值，
只是把它们看成 $u$ 和 $t$ 的函数而不是 $r$ 和 $t$ 的函数。
所以式子里仍然写 $C$，只是含义变成 $C(u,t)$。

\subsection{第五步：边界条件也要改写}

表面在 $r=R$ 处，对应 $u=1$。把 $u=1$ 代入式 \eqref{eq:drdx}，

\begin{equation}
\left.\frac{\partial}{\partial r}\right|_{u=1}
= \frac{2\sqrt{1}}{R}\left.\frac{\partial}{\partial u}\right|_{u=1}
= \frac{2}{R}\left.\frac{\partial}{\partial u}\right|_{u=1}.
\end{equation}

温度的表面条件原来是

\begin{equation}
-k\left.\frac{\partial T}{\partial r}\right|_{r=R}
= h\left(T(R,t)-T_{\mathrm{air}}(t)\right).
\end{equation}

把左边换成 $u$ 的导数，

\begin{equation}
-k\cdot\frac{2}{R}\left.\frac{\partial T}{\partial u}\right|_{u=1}
= h\left(T_1-T_{\mathrm{air}}(t)\right),
\end{equation}

其中 $T_1$ 表示 $u=1$ 处的温度，也就是表面温度。两边同乘 $-R/2$，

\begin{equation}
k\left.\frac{\partial T}{\partial u}\right|_{u=1}
= -\frac{hR}{2}\left(T_1-T_{\mathrm{air}}(t)\right)
= \frac{hR}{2}\left(T_{\mathrm{air}}(t)-T_1\right).
\label{eq:bcu-T}
\end{equation}

水分条件同样处理，

\begin{equation}
D\left.\frac{\partial C}{\partial u}\right|_{u=1}
= \frac{h_mR}{2}\left(C_{\mathrm{air}}(t)-C_1\right).
\label{eq:bcu-C}
\end{equation}

写成这种形式以后，左右两边都不含根号和分数，
程序里直接照抄这两个式子就行。

\subsection{第六步：圆心处还需要加条件吗}

不需要。

原因在于，当 $u=0$ 时，式 \eqref{eq:mass-u} 右边的通量因子
$uD\dfrac{\partial C}{\partial u}$ 里有一个 $u$，等于零。
这意味着从圆心流出的通量天然就是零，
正好就是对称条件 $\dfrac{\partial C}{\partial r}\big|_{r=0}=0$ 的含义。

换句话说，对称条件已经自动包含在方程里了，不需要另外写。

\keypoint
坐标变换做了三件事。
第一，把 $1/r$ 消掉了，圆心不再是特殊点。
第二，两个方程形式完全一样，只是系数不同，程序可以共用一套代码。
第三，表面条件化成整洁的代数式。

% =====================================================================
\section{谱方法：用很少的点算得很准}

有了新坐标，接下来要把区间 $0\le u\le1$ 切成有限个点，
在点上求近似解。这一章讲怎么切、怎么把导数变成代数运算。

\subsection{基本想法：用多项式代替未知函数}

假设我们不知道 $C(u,t)$ 的具体形式，但知道它在若干个点上的值。
能不能构造一个多项式，让它在这几个点上的值正好等于已知值，
然后用这个多项式代表未知函数？

可以。这个多项式叫插值多项式，这几个点叫插值节点。

一旦有了多项式，求导数就变成对多项式求导，结果仍然是多项式，
可以精确算出它在各点的值。这样微分运算就变成了有限的代数运算。

\subsection{节点怎么选：一个反面教材}

最自然的想法是把点均匀分布。比如在 $0$ 到 $1$ 之间取 $11$ 个等距点。
这种做法在点数少的时候还行，点数一多就出问题。

问题出在函数形状复杂的地方。均匀分布的点无法同时照顾区间两端和中间，
当多项式次数升高时，两端会出现剧烈的上下摆动。
这个现象叫龙格现象，是插值理论里的经典结论。

\why{为什么会摆动}
直觉的解释是：区间两端的点比较少的信息去约束多项式，
而高次多项式在端点的变化率很大，
稍微改变一点系数就会让端点附近剧烈起伏。
均匀分布没有在两端加密，所以控制不住。

\subsection{切比雪夫节点}

解决办法是在区间两端加密。具体做法是取

\begin{equation}
x_j = \cos\frac{j\pi}{N}, \qquad j=0,1,2,\dots,N .
\label{eq:chebnode}
\end{equation}

这些点分布在 $-1$ 到 $1$ 之间。当 $j=0$ 时 $x_0=\cos0=1$，
当 $j=N$ 时 $x_N=\cos\pi=-1$。

这些点有一个特点：靠近两端时，$\cos$ 的变化变慢，
所以点会自然地在两端聚集。可以验证，最靠近端点的两个点之间的距离
大约是 $\dfrac{\pi^2}{2N^2}$，而中间的间距大约是 $\dfrac{\pi}{N}$。
两端比中间密了约 $N$ 倍。

再把它们映射到区间 $0$ 到 $1$：

\begin{equation}
u_j = \frac{1-x_j}{2}.
\label{eq:map01}
\end{equation}

当 $x=1$ 时 $u=0$，当 $x=-1$ 时 $u=1$。这样 $u_0=0$，$u_N=1$，
正好对应圆心和表面。

\why{映射公式里的 1 和 2 是怎么来的}
把区间 $-1$ 到 $1$ 线性映射到 $0$ 到 $1$，
需要先加 $1$ 把区间移到 $0$ 到 $2$，再除以 $2$ 缩放到 $0$ 到 $1$。
也就是 $u=\dfrac{x+1}{2}$。
而我们的 $x_j=\cos$ 是从 $1$ 递减到 $-1$ 的，
所以写成 $u=\dfrac{1-x}{2}$ 正好让 $u$ 从 $0$ 递增到 $1$。

\subsection{插值多项式的唯一性}

有一个重要的事实：给定 $N+1$ 个互不相同的节点和对应的函数值，
次数不超过 $N$ 的插值多项式是唯一的。

理由是这样。设两个不同的多项式都满足条件，它们的差也是一个次数不超过 $N$
的多项式，但在这个 $N+1$ 个节点上都等于零。
一个次数不超过 $N$ 的非零多项式最多只有 $N$ 个根，
现在有了 $N+1$ 个根，只能是恒等于零的多项式。所以两个多项式其实是同一个。

这个事实保证了插值结果是确定的，不受计算方法影响。

\subsection{拉格朗日插值公式}

插值多项式可以写成

\begin{equation}
p(u) = \sum_{j=0}^{N} C_j\,\ell_j(u),
\label{eq:lagrange}
\end{equation}

其中 $C_j$ 是函数在第 $j$ 个节点上的值，$\ell_j(u)$ 叫基函数，定义是

\begin{equation}
\ell_j(u) = \prod_{\substack{k=0\\k\ne j}}^{N}\frac{u-u_k}{u_j-u_k}.
\label{eq:basis}
\end{equation}

这个式子的意思需要解释一下。

分子是 $N$ 个因子的乘积，每个因子是 $u$ 减去一个别的节点的位置。
当 $u$ 取第 $j$ 个节点时，分子里出现因子 $u_j-u_j=0$，整个分子等于零。

分母是把分子里的 $u$ 换成 $u_j$ 得到的常数。

所以当 $u=u_j$ 时，分子分母相同，$\ell_j(u_j)=1$。
而当 $u=u_k$（$k$ 不等于 $j$）时，分子里有一个因子 $u_k-u_k=0$，
所以 $\ell_j(u_k)=0$。

把这两条合起来，代入式 \eqref{eq:lagrange}，当 $u=u_j$ 时，

\begin{equation}
p(u_j) = C_j\cdot1 + \sum_{k\ne j}C_k\cdot0 = C_j .
\end{equation}

插值条件得到满足。

\subsection{从插值多项式得到导数}

对式 \eqref{eq:lagrange} 求导，

\begin{equation}
p'(u) = \sum_{j=0}^{N}C_j\,\ell_j'(u).
\end{equation}

在节点 $u_i$ 处取值，

\begin{equation}
p'(u_i) = \sum_{j=0}^{N}C_j\,\ell_j'(u_i).
\end{equation}

把 $\ell_j'(u_i)$ 记作 $D_{ij}$，就得到一个矩阵形式的关系：

\begin{equation}
\begin{pmatrix} p'(u_0)\\ p'(u_1)\\ \vdots\\ p'(u_N)\end{pmatrix}
=
\begin{pmatrix}
D_{00} & D_{01} & \cdots & D_{0N}\\
D_{10} & D_{11} & \cdots & D_{1N}\\
\vdots & \vdots & \ddots & \vdots\\
D_{N0} & D_{N1} & \cdots & D_{NN}
\end{pmatrix}
\begin{pmatrix} C_0\\ C_1\\ \vdots\\ C_N\end{pmatrix}.
\label{eq:dmatrix}
\end{equation}

这个矩阵叫微分矩阵。它的作用是：输入一组函数值，输出一组导数值。
求导这个微分运算，被彻底变成了矩阵乘法。

\subsection{微分矩阵的元素公式}

对切比雪夫节点，微分矩阵的元素有封闭形式的公式。记

\begin{equation}
c_0 = 2,\qquad c_N = 2,\qquad c_j = 1 \quad (1\le j\le N-1).
\end{equation}

当 $i$ 不等于 $j$ 时，

\begin{equation}
D_{ij} = \frac{c_i}{c_j}\cdot\frac{(-1)^{i+j}}{x_i-x_j}.
\label{eq:dij}
\end{equation}

当 $i$ 等于 $j$ 时，不能直接用这个公式，因为分母为零。
对角线元素用另一条性质确定：常数函数的导数处处为零。
把常数函数 $f\equiv1$ 代入式 \eqref{eq:dmatrix}，
左边全是零，右边是对每一行求和，于是

\begin{equation}
\sum_{j=0}^{N}D_{ij} = 0 \quad\text{对每个 }i,
\end{equation}

所以

\begin{equation}
D_{ii} = -\sum_{\substack{j=0\\j\ne i}}^{N}D_{ij}.
\label{eq:dii}
\end{equation}

\why{为什么要用这个办法求对角元}
因为直接计算对角元需要求极限，数值上不稳定。
用行和为零这条性质，既精确又简单。
这条性质来自一个简单事实：常数的导数为零，
而插值多项式对常数函数是精确的。

\subsection{记号的换算}

式 \eqref{eq:dij} 里的 $x$ 是切比雪夫变量。我们的自变量是 $u$，
两者关系是 $u=\dfrac{1-x}{2}$，反过来 $x=1-2u$。

对 $u$ 求导和对 $x$ 求导的关系由链式法则给出：

\begin{equation}
\frac{\partial}{\partial u}
= \frac{\partial x}{\partial u}\cdot\frac{\partial}{\partial x}
= \frac{\partial(1-2u)}{\partial u}\cdot\frac{\partial}{\partial x}
= -2\frac{\partial}{\partial x}.
\end{equation}

所以

\begin{equation}
D_u = -2D_x .
\label{eq:du}
\end{equation}

程序里先算 $D_x$，再乘 $-2$ 得到 $D_u$，最后对 $u$ 求导就用 $D_u$。

\subsection{一个三节点的手算例子}

上面的公式看着抽象。这一段把节点数取成 $N=2$，
也就是只用三个点，把整个过程从零算一遍。
每一个数字都会写出来，你可以拿计算器逐个复核。

\textbf{第一步，算节点位置。}

按式 \eqref{eq:chebnode}，让 $j$ 依次取 $0$、$1$、$2$：

\begin{equation}
x_0=\cos\frac{0\cdot\pi}{2}=\cos0=1,
\end{equation}

\begin{equation}
x_1=\cos\frac{1\cdot\pi}{2}=\cos\frac{\pi}{2}=0,
\end{equation}

\begin{equation}
x_2=\cos\frac{2\cdot\pi}{2}=\cos\pi=-1 .
\end{equation}

再按式 \eqref{eq:map01} 换算到 $u$ 坐标：

\begin{equation}
u_0=\frac{1-1}{2}=0,\qquad
u_1=\frac{1-0}{2}=0.5,\qquad
u_2=\frac{1-(-1)}{2}=1 .
\end{equation}

三个节点分别落在圆心、径向中点、表面。

\textbf{第二步，算系数 $c$。}

按定义，首尾两个取 $2$，中间取 $1$，所以

\begin{equation}
c_0=2,\qquad c_1=1,\qquad c_2=2 .
\end{equation}

\textbf{第三步，算节点之间的差值 $x_i-x_j$。}

把所有九个组合列成表。先写对角线为零，稍后修正。

\begin{center}
\begin{tabular}{@{}c|ccc@{}}
\toprule
$i$ 与 $j$ & $j=0$ & $j=1$ & $j=2$ \\
\midrule
$i=0$ & $0$ & $1-0=1$ & $1-(-1)=2$ \\
$i=1$ & $0-1=-1$ & $0$ & $0-(-1)=1$ \\
$i=2$ & $-1-1=-2$ & $-1-0=-1$ & $0$ \\
\bottomrule
\end{tabular}
\end{center}

\textbf{第四步，算非对角线上的元素。}

按式 \eqref{eq:dij}，公式是

\begin{equation}
D_{ij}=\frac{c_i}{c_j}\cdot\frac{(-1)^{i+j}}{x_i-x_j}.
\end{equation}

逐个算六个非对角线元素。

$i=0$，$j=1$：

\begin{equation}
D_{01}=\frac{c_0}{c_1}\cdot\frac{(-1)^{1}}{x_0-x_1}
=\frac{2}{1}\cdot\frac{-1}{1}=-2 .
\end{equation}

$i=0$，$j=2$：

\begin{equation}
D_{02}=\frac{2}{2}\cdot\frac{(-1)^{2}}{2}=1\times\frac{1}{2}=0.5 .
\end{equation}

$i=1$，$j=0$：

\begin{equation}
D_{10}=\frac{1}{2}\cdot\frac{(-1)^{1}}{-1}=0.5\times\frac{-1}{-1}=0.5 .
\end{equation}

$i=1$，$j=2$：

\begin{equation}
D_{12}=\frac{1}{2}\cdot\frac{(-1)^{3}}{1}=0.5\times(-1)=-0.5 .
\end{equation}

$i=2$，$j=0$：

\begin{equation}
D_{20}=\frac{2}{2}\cdot\frac{(-1)^{2}}{-2}=1\times\frac{1}{-2}=-0.5 .
\end{equation}

$i=2$，$j=1$：

\begin{equation}
D_{21}=\frac{2}{1}\cdot\frac{(-1)^{3}}{-1}=2\times\frac{-1}{-1}=2 .
\end{equation}

\why{程序里为什么不写这个公式}
程序里写的是先把 $(-1)^j$ 乘进系数 $c$，
于是 $c_i/c_j$ 这一个因子就同时包含了 $\dfrac{c_i}{c_j}$ 和 $(-1)^{i+j}$，
公式变成一次除法。这样写更短，结果完全一样。
如果用上面这个公式逐个手算，得到的是同一组数字。

\textbf{第五步，用行和为零求对角线元素。}

对角线上不能直接用上面的公式，因为分母为零。
改用另一条性质：常数函数的导数处处为零。
把常数函数 $f\equiv1$ 代入，左边每行都是零，
右边是矩阵每一行的和，于是每一行的和必须等于零：

\begin{equation}
D_{ii}=-\sum_{j\ne i}D_{ij}.
\end{equation}

第一行：

\begin{equation}
D_{00}=-(D_{01}+D_{02})=-(-2+0.5)=-(-1.5)=1.5 .
\end{equation}

第二行：

\begin{equation}
D_{11}=-(D_{10}+D_{12})=-(0.5-0.5)=-0=0 .
\end{equation}

第三行：

\begin{equation}
D_{22}=-(D_{20}+D_{21})=-(-0.5+2)=-(1.5)=-1.5 .
\end{equation}

把全部元素拼起来，得到

\begin{equation}
D_x=
\begin{pmatrix}
1.5 & -2 & 0.5\\
0.5 & 0 & -0.5\\
-0.5 & 2 & -1.5
\end{pmatrix}.
\end{equation}

\textbf{第六步，验证这个矩阵对不对。}

用两个已知导数的函数来检验。

第一个用常数函数 $f(x)=1$，它的导数是零。
把 $f$ 在三个节点上的值写成列向量 $(1,1,1)$，用矩阵去乘：

\begin{equation}
\begin{pmatrix}1.5 & -2 & 0.5\\ 0.5 & 0 & -0.5\\ -0.5 & 2 & -1.5\end{pmatrix}
\begin{pmatrix}1\\1\\1\end{pmatrix}
=
\begin{pmatrix}1.5-2+0.5\\ 0.5+0-0.5\\ -0.5+2-1.5\end{pmatrix}
=
\begin{pmatrix}0\\0\\0\end{pmatrix}.
\end{equation}

三个分量都是零，正确。

第二个用一次函数 $f(x)=x$，它的导数是常数 $1$。
$f$ 在三个节点上的值是 $(1,0,-1)$，

\begin{equation}
\begin{pmatrix}1.5 & -2 & 0.5\\ 0.5 & 0 & -0.5\\ -0.5 & 2 & -1.5\end{pmatrix}
\begin{pmatrix}1\\0\\-1\end{pmatrix}
=
\begin{pmatrix}1.5+0-0.5\\ 0.5+0+0.5\\ -0.5+0+1.5\end{pmatrix}
=
\begin{pmatrix}1\\1\\1\end{pmatrix}.
\end{equation}

三个分量都是 $1$，也正确。

\textbf{第七步，换成对 $u$ 求导的矩阵。}

按式 \eqref{eq:du}，把上面每个元素乘 $-2$：

\begin{equation}
D_u=-2D_x=
\begin{pmatrix}
-3 & 4 & -1\\
-1 & 0 & 1\\
1 & -4 & 3
\end{pmatrix}.
\end{equation}

再用一次函数检验。在 $u$ 坐标下取 $f(u)=u$，
它在三个节点上的值是 $(0,\ 0.5,\ 1)$，而 $\dfrac{\mathrm{d}u}{\mathrm{d}u}=1$，

\begin{equation}
\begin{pmatrix}-3 & 4 & -1\\ -1 & 0 & 1\\ 1 & -4 & 3\end{pmatrix}
\begin{pmatrix}0\\0.5\\1\end{pmatrix}
=
\begin{pmatrix}0+2-1\\ 0+0+1\\ 0-2+3\end{pmatrix}
=
\begin{pmatrix}1\\1\\1\end{pmatrix}.
\end{equation}

三个分量都是 $1$，说明程序里那行 $D_u=-2D_x$ 得到的结果是对的。

\keypoint
三个节点的手算结果给出两条检验标准。
第一，微分矩阵每一行的和必须等于零，因为常数的导数为零。
第二，把节点坐标本身当作函数代进去，结果必须全是 $1$，因为一次函数的导数是常数。
这两条只要有一条不满足，矩阵就一定写错了。

\subsection{为什么谱方法精度高}

常规的差分方法，比如用相邻两点做差商近似导数，
误差是步长的平方量级。步长减半，误差降到四分之一。
要想精度提高一千倍，得把点数增加三十多倍。

谱方法不一样。如果函数足够光滑，
插值多项式的误差随节点数增加而指数下降。
也就是说，点数稍微增加一点，误差就急剧变小。

原因可以这样理解。差分方法只用了相邻几个点的信息，
而谱方法的插值多项式用了所有点的信息，
等价于用整个区间上的信息去估计每一点的导数，
所以效率高得多。

\why{本题能用谱方法吗}
可以，而且效果很好。原因在第五章已经埋下伏笔。
在 $u$ 坐标下，解是关于 $u$ 的解析函数，
具体说，第一类贝塞尔函数 $J_0$ 的展开式

\begin{equation}
J_0\!\left(\frac{\beta r}{R}\right)
= \sum_{k\ge0}\frac{(-1)^k}{k!\,k!}\left(\frac{\beta^2}{4}\right)^{k}u^{k}
\end{equation}

只含 $u$ 的整数次幂，是一个处处收敛的幂级数。
这类函数正是多项式插值最擅长逼近的对象。
实际计算中，取 $N=8$，也就是只用 $9$ 个点，
与精确解的偏差就已经降到 $10^{-10}$ 的量级。

\keypoint
谱方法的三步是：选切比雪夫节点，构造微分矩阵，用矩阵乘法代替求导。
本题里 $N=8$ 就达到 $10^{-10}$ 的精度，
而常规方法要达到同样精度需要上千个点。

% =====================================================================
\section{时间推进：从当前时刻推到下一时刻}

到上一章为止，我们把空间方向的导数处理完了。
剩下时间方向的导数 $\dfrac{\partial C}{\partial t}$ 还没处理。
这一章讲怎么处理。

\subsection{半离散化}

先只对空间做离散。在 $N+1$ 个节点上，方程

\begin{equation}
\frac{\partial C}{\partial t}=\frac{4}{R^2}\frac{\partial}{\partial u}\left(uD\frac{\partial C}{\partial u}\right)
\end{equation}

变成 $N+1$ 个常微分方程。比如第 $j$ 个节点上的方程是

\begin{equation}
\frac{\mathrm{d}C_j}{\mathrm{d}t}
= \frac{4}{R^2}\left[\text{第 }j\text{ 个节点上的空间算子离散值}\right].
\label{eq:semi}
\end{equation}

这个过程叫半离散化，意思是空间已经离散，时间还是连续的。

\why{为什么先做空间再做时间}
因为空间离散化以后，问题的规模就定下来了，
只剩下一组关于时间的常微分方程。
这种方程有很多成熟的求解程序可以直接用，
不需要自己从零开始写。这也是一种分工策略：
难的部分自己做，成熟的部分用现成的。

\subsection{最简单的方法：显式欧拉法}

把时间也切成小段，步长记作 $\Delta t$。
最简单的方法是用前向差商近似时间导数：

\begin{equation}
\frac{\mathrm{d}C_j}{\mathrm{d}t}\approx\frac{C_j^{n+1}-C_j^{n}}{\Delta t},
\label{eq:forward}
\end{equation}

其中上标 $n$ 表示第 $n$ 个时刻，$n+1$ 表示下一个时刻。

把式 \eqref{eq:forward} 代入式 \eqref{eq:semi}，得到

\begin{equation}
C_j^{n+1} = C_j^{n} + \Delta t\cdot
\frac{4}{R^2}\left[\text{空间算子}（C^n）\right].
\label{eq:euler}
\end{equation}

右边全部是已知的旧值，直接算出来就是新值，不需要解方程。
这个方法叫显式欧拉法。

\subsection{显式欧拉法的问题：稳定性}

显式欧拉法有一个致命缺点。它可以写成

\begin{equation}
C^{n+1} = \left(I+\Delta t\,L\right)C^{n},
\end{equation}

其中 $L$ 是空间离散矩阵，$I$ 是单位矩阵。
要让解不发散，要求 $\Delta t$ 足够小，
具体来说要求 $\Delta t$ 小于某个与 $\Delta u^2$ 同量级的阈值。

而谱方法的节点在两端非常密集，最小的节点间距大约是 $\dfrac{\pi^2}{2N^2}$，
它对应的稳定步长阈值正比于这个量的平方，
也就是正比于 $\dfrac{1}{N^4}$。

算一下量级。取 $N=64$，$\dfrac{1}{N^4}$ 约等于 $6\times10^{-8}$。
在量纲上还要乘以 $\dfrac{R^2}{4D}$，大约是 $2000$ 秒。
两者相乘得到约 $1.2\times10^{-4}$ 秒。
也就是说每一秒要算将近一万步，而整个过程要算 $20$ 万秒。
这个计算量是无法接受的。

这种要求时间步长远小于物理过程时间尺度的现象，叫做刚性。
本题就是典型的刚性问题。

\why{刚性的物理来源}
表面处的干燥非常快，因为一开始表面含水率是 $2.55$，
而空气中的含湿量只有 $0.0196$，落差很大，表面迅速失水。
在很短的时间内，表面附近形成一层很薄的水分变化区。
要准确刻画这一层，时间步长必须很小。
但整个干燥过程长达五十多个小时。
快慢两个尺度相差极大，这就是刚性的根源。

\subsection{隐式方法}

换一个思路。用后向差商近似时间导数：

\begin{equation}
\frac{\mathrm{d}C_j}{\mathrm{d}t}\approx\frac{C_j^{n+1}-C_j^{n}}{\Delta t}.
\end{equation}

式子的形式看起来和式 \eqref{eq:forward} 一样，区别在于右端空间算子取的是哪个时刻的值。
显式方法取旧时刻 $n$，隐式方法取新时刻 $n+1$：

\begin{equation}
\frac{C_j^{n+1}-C_j^{n}}{\Delta t}
= \frac{4}{R^2}\left[\text{空间算子}（C^{n+1}）\right].
\label{eq:implicit}
\end{equation}

现在右边含有未知的新值，两边都有未知量，
需要解一个方程组才能得到新值。这就是隐式的含义。

\why{隐式方法为什么能允许大得多的步长}
可以这样理解。显式方法相当于说我用现在的趋势直接外推，
如果趋势变化很快，外推就会跑偏。
隐式方法相当于说我先猜一个新值，然后检查它是否满足方程，
不满足就调整。这种自我修正的机制让方法对快速变化有更强的抵抗力。
数学上可以证明，后向欧拉法是无条件稳定的，
不管步长多大解都不会发散。

代价是每一步都要解方程组。但本题的方程组规模很小，
用现成的线性代数程序在毫秒级就能解完，完全可以接受。

\subsection{更高阶的方法：BDF}

后向欧拉法只有一阶精度，也就是说误差与步长成正比。
步长减半，误差减半。要提高精度，可以用更高阶的公式。

BDF 是向后差分公式的缩写。它的想法是利用前面若干个时刻的解，
构造一个更高精度的差商近似。

一阶 BDF 就是后向欧拉法，用到 $1$ 个历史时刻。
二阶 BDF 用到 $2$ 个历史时刻，公式是

\begin{equation}
\frac{3C^{n+1}-4C^{n}+C^{n-1}}{2\Delta t}
= \frac{4}{R^2}\left[\text{空间算子}（C^{n+1}）\right].
\end{equation}

分子上的系数 $3$、$-4$、$1$ 来自对二次插值多项式求导，
不是随便凑出来的。三阶、四阶、五阶的形式类似，
分别用到 $3$、$4$、$5$ 个历史时刻。

阶数越高，精度越高，但稳定性区域越小，对步长的限制越多。
实践中的做法是在程序里同时实现多个阶数，
根据当前的误差估计自动选择用几阶。

\subsection{自适应步长}

步长固定不变有一个弊端。干燥刚开始时变化剧烈，需要小步长；
过了两个小时，温度和水分都变化很慢，大步长就够用了。
如果全程都用小步长，会浪费大量时间。

自适应步长的做法是：先按当前步长算一步，
同时估算这一步的误差。如果误差超过容忍值，就把步长减小重算；
如果误差远小于容忍值，就把步长增大。

误差估计的常用办法是比较两个不同阶数方法的结果，
两者的差可以作为误差的近似。这个技巧叫嵌入式方法。

程序里的两个控制参数是

\begin{equation}
\mathrm{rtol}\ \text{（相对容差）}, \qquad \mathrm{atol}\ \text{（绝对容差）}.
\end{equation}

判据是每一步的误差要满足

\begin{equation}
\text{误差} \le \mathrm{atol} + \mathrm{rtol}\times\left|\text{解}\right|.
\end{equation}

本题取的数值是 $\mathrm{rtol}=10^{-11}$，
温度分量的 $\mathrm{atol}$ 取 $10^{-11}$，
水分分量的 $\mathrm{atol}$ 取 $10^{-13}$。

\why{为什么水分和温度的容差取得不一样}
因为两者的量级不一样。温度是几十摄氏度，
水分浓度是零点几到两点几。如果都用同一个绝对容差，
水分那边可能过于宽松，温度那边可能过于严格。
按各自量级分别设置更合理。

\subsection{一个必须注意的细节：系数用什么时刻的值}

方程里的 $D$、$k$、$\rho c_p$ 都依赖 $C$ 和 $T$，
而 $C$ 和 $T$ 本身是未知的。这就带来一个选择：

可以在每个时间步开始时，用旧的 $C$ 和 $T$ 算出系数，
然后在整个时间步内把系数当作常数使用。这种做法叫系数冻结。
它的优点是每一步只需要解一次线性方程组。

也可以在每一步内部迭代若干次，每次都用最新的 $C$ 和 $T$ 重新计算系数。
这种做法更准，但每一步要解若干次方程组。

本题的主求解器采用的是一种折中：右端项每次被调用时，
都用当时传入的 $C$ 和 $T$ 重新计算系数。
由于自适应积分器在每一步内部会多次调用右端项，
实际上相当于做了隐式的迭代，精度更高。
代价是每一步的计算量稍大，但程序更简洁，不容易出错。

\keypoint
时间推进的核心结论是：本题是刚性问题，
必须用隐式方法；用自适应变阶的 BDF 方法可以在保证精度的同时节省计算量。

% =====================================================================
\section{烘干时间怎么确定}

前几章解决的是怎么算出任意时刻的温度和水分场。
问题三和问题四要回答的是另一个问题：
什么时候算烘干完成。这一章讲这个判断怎么做。

\subsection{判据的数学表述}

题目说药材各处的水分浓度应低于 $0.15$ 千克每千克。
这句话要翻译成数学语言。

关键字是各处。也就是整个径向范围内每一个位置都要满足，
不能只看平均值，也不能只看表面。

写成式子就是

\begin{equation}
\max_{0\le r\le R_0}C(r,t) < 0.15 .
\end{equation}

烘干时间是第一次满足这个条件的那一时刻，写成

\begin{equation}
t_{\mathrm{dry}}=\min\left\{t>0:\ \max_{0\le r\le R_0}C(r,t)<0.15\right\}.
\label{eq:tdry}
\end{equation}

这个式子的读法是：在所有使最大值小于 $0.15$ 的时刻里，取最小的那一个。

\why{为什么必须逐点判断而不能用平均值}
因为干燥是从表面往内部推进的。
在干燥的中后期，表面已经很干，而中心还比较湿。
如果看平均值，平均值可能早就低于 $0.15$ 了，
但中心还没达标，药材实际上还是湿的。

本题的实际数据是：截面平均含水率在 $35.2$ 小时就降到了 $0.15$，
而中心要到 $57.1$ 小时才达标，晚了将近 $22$ 小时。
用平均值判断会把时间低估约 $38\%$。

\subsection{哪一点最后干}

从物理上可以预判：中心应该是最后干的地方。
理由是水分要从内部扩散到表面才能被带走。
中心到表面的距离最长，扩散路程最长，
而且它周围都是含水率更高的地方，浓度梯度最小，驱动力最弱。

但是本文没有把这个预判直接写进程序，
而是在每一个时刻对全部节点取最大值，看最大值出现在哪里。
这样做的原因是：如果预判错了，程序不会跟着错。
实际运行的每一时刻结果都显示最大值出现在中心，
预判得到了验证。

\why{为什么要在全部节点上取最大而不是只在输出位置上取最大}
因为输出的位置是每隔 $0.1$ 厘米一个点，共 $21$ 个点。
而计算的节点有 $N+1$ 个，在两端更密集。
如果只在 $21$ 个输出点上取最大，
有可能真正的最大值出现在两个输出点之间而没有被发现。
虽然本题中最大值总在中心，但把范围取全是一个更安全的做法。

\subsection{怎么求这个时刻}

式 \eqref{eq:tdry} 定义的是一个事件的时刻。
程序里的处理办法是构造一个函数

\begin{equation}
e(t) = \max_{0\le j\le N}C_j(t) - 0.15 .
\end{equation}

当 $e(t)>0$ 时说明还没干，当 $e(t)<0$ 时说明已经干了。
烘干时间就是 $e(t)$ 第一次从正变负的时刻，
也就是 $e(t)=0$ 的第一个根。

求根用二分法或者布伦特方法。这两种方法都是标准的数值方法，
程序里由积分器自动调用。

两个方法的基本原理很简单。二分法是：
如果知道 $e$ 在 $t_1$ 处为正、在 $t_2$ 处为负，
那么根一定在两者之间。取中点 $t_3$，
看 $e(t_3)$ 的符号，就能把根所在的范围缩小一半。
反复这样做，区间会越来越小，最后得到一个足够精确的近似值。

\why{为什么不用在网格序列上做外推}
有一种常见做法是：用几组不同的网格算若干次，
得到若干个烘干时间，然后按某种规律外推到一个无限细的网格。
这种做法需要假设收敛的阶数，如果假设错了结果就会偏。

本文的做法更直接：直接用自适应积分器给出的高精度解来求根，
不需要任何关于收敛阶的假设。
后来在第四部分会看到，这两种做法的结果确实相差约 $15$ 秒，
差异来自外推所用的序列本身没有收敛。

\keypoint
烘干时间的判定要注意三点。
第一，判据是逐点的最大值，不是平均值。
第二，程序在所有节点上取最大，不预设极值位置。
第三，用求根而不是外推来确定时刻。

\subsection{第二部分小结}

第二部分解决了从方程到程序的全部原理问题：

\begin{center}
\begin{tabular}{@{}lll@{}}
\toprule
问题 & 办法 & 在本文中的名称 \\
\midrule
圆心处 $1/r$ 的麻烦 & 换元 $u=(r/R)^2$ & 坐标正则化 \\
空间导数怎么算 & 切比雪夫节点加微分矩阵 & 谱配置法 \\
时间方向怎么推进 & 隐式变阶 BDF 加自适应步长 & 自适应时间积分 \\
导数怎么取近似 & 通量散度形式 & 守恒型离散 \\
烘干时刻怎么确定 & 构造函数求零点 & 终端事件求交 \\
\bottomrule
\end{tabular}
\end{center}

% =====================================================================
\section{问题四的收缩怎么处理}

问题一到问题三都假设药材半径不变。问题四不一样，
题目给了附件二，里面记录了半径随时间的变化，
从最初的 $2.000$ 厘米一直降到 $1.198$ 厘米，缩了四成。

半径变化意味着求解区域的边界在移动。
这一章专门讲怎么处理这种情况。

\subsection{第一种想法为什么是错的}

最自然的想法是：半径坐标下的方程不是已经写好了吗，
直接让 $R$ 随时间变化不就行了。

也就是说，直接使用式 \eqref{eq:mass}

\begin{equation}
\frac{\partial C}{\partial t}
= \frac{1}{r}\frac{\partial}{\partial r}\left(r\,D\,\frac{\partial C}{\partial r}\right)
\end{equation}

再把表面条件里的 $R$ 换成 $R(t)$。

这个做法是错的。原因要讲清楚，因为直觉上很难发现。

式 \eqref{eq:mass} 里的时间导数是在固定空间位置求的。
意思是：站在空间里某个固定的点，看那里的含水率怎么变。

但药材在收缩，干物质在移动。
水分是跟着干物质一起走的。
站在固定的空间点上看，含水率的变化有两个来源：
一个是水分的扩散，另一个是干物质带着水分从旁边走过。

后一个来源在原方程里没有被描述，因为原方程推导时假设了干物质不动。

\subsection{用一个极端例子把这个错误暴露出来}

为了看清楚，考虑一个极端情况：假设扩散系数等于零。

此时水分完全不能扩散，全部被困在干物质里面。
那么每个材料点的含水率永远不会变。

再假设初始含水率处处相同。这符合题目条件，题目给的初值就是
$C_0=2.55$ 在整根药材上处处一样。

两个条件合起来得到什么结论？每个材料点的含水率都一直是 $2.55$，
而且所有材料点都一样。
换句话说，含水率的分布从头到尾都是一个常数。

既然分布是常数，那么在任何时刻、任何位置，
含水率都不应该发生变化，也就是

\begin{equation}
\left.\frac{\partial C}{\partial t}\right|_r = 0 .
\end{equation}

现在看错误的方程给出什么。把 $D=0$ 代入式 \eqref{eq:mass}，右边变成零，
所以方程给出 $\dfrac{\partial C}{\partial t}=0$。

咦，这里看起来是对的。问题出在哪里？

问题出在上面这个推理本身有漏洞。
我们说含水率分布是常数，指的是在材料坐标系下的分布。
在固定空间坐标下，情况并不是这样。

具体说，本来在半径 $1.5$ 厘米处的干物质，随着收缩慢慢移到半径 $1.0$ 厘米处。
它带着自己的含水率一起走。
如果初始含水率不均匀，比如中心是 $2.55$、表面是 $1.5$，
那么在固定的空间点 $r=1.0$ 厘米处，
一开始那里是内部的水（含水率高），后来流过来的是表层的水（含水率低）。
含水率自然就变了。

所以固定空间坐标下的时间导数不等于零，
错误的方程给出的零是错的。

\why{这个错误为什么不容易发现}
因为当初始含水率处处相同时，两种写法给出的结果都是常数，
看起来都合理。
只有在含水率不均匀时错误才会暴露。
而实际干燥过程中含水率恰恰是不均匀的，
所以错误的模型会给出错误的干燥时间。

本文的测试结果是：如果误在固定半径坐标下直接写方程，
算出的烘干时间会明显偏长，偏差达到数倍。

\subsection{正确的做法：跟着干物质走}

既然干物质在动，那就把坐标系挂在干物质上。

引入物质坐标

\begin{equation}
\xi = \frac{r}{R(t)} .
\end{equation}

这个坐标的含义是相对位置。当 $r=0$ 时 $\xi=0$，当 $r=R(t)$ 时 $\xi=1$。
不管药材怎么缩，圆心永远是 $\xi=0$，表面永远是 $\xi=1$。

关键性质是：对同一个材料点，$\xi$ 永远不变。

验证一下。设某个材料点在 $t=0$ 时位于半径 $r_0$ 处，
那么它的标签是 $\xi_0=r_0/R_0$。
由于收缩是各向同性的，也就是所有位置按同样的比例缩小，
这个点在时刻 $t$ 的位置是

\begin{equation}
r(t) = r_0\cdot\frac{R(t)}{R_0} .
\end{equation}

代入 $\xi$ 的定义，

\begin{equation}
\xi = \frac{r(t)}{R(t)} = \frac{r_0 R(t)/R_0}{R(t)} = \frac{r_0}{R_0} = \xi_0 .
\end{equation}

分子分母里的 $R(t)$ 约掉了，结果与时间无关。
所以 $\xi$ 确实是一个不随时间变化的标签。

\why{各向同性这个假设为什么必须写明}
如果药材在长度方向不缩、只在半径方向缩，
那么材料点的运动就不是简单的按比例缩小，
$r(t)$ 和 $R(t)$ 的比值就不再是常数，
上面的推导就不成立，方程会多出额外的项。
题目只给了半径方向的收缩数据，
按各向同性处理是同口径的做法，论文里明确写成了假设。

\subsection{干物质守恒给出的一个重要关系}

现在对一个材料壳层做干物质衡算。

取标签在 $\xi$ 到 $\xi+\mathrm{d}\xi$ 之间的那一层材料。
它在当前构形下的体积（单位长度）是多少？

径向位置是 $r=\xi R$，厚度是 $\mathrm{d}r = R\,\mathrm{d}\xi$。
体积（单位长度）等于周长乘以厚度，

\begin{equation}
\mathrm{d}V = 2\pi r\,\mathrm{d}r = 2\pi\,\xi R\cdot R\,\mathrm{d}\xi
= 2\pi R^2\xi\,\mathrm{d}\xi .
\end{equation}

这个壳层里的干物质质量是

\begin{equation}
\mathrm{d}m_d = \rho_d\,\mathrm{d}V = \rho_d\,2\pi R^2\xi\,\mathrm{d}\xi .
\end{equation}

干物质不会迁移，也不会凭空产生或消失，
所以这个质量不随时间变化：

\begin{equation}
\frac{\partial}{\partial t}\left(\rho_d\,2\pi R^2\xi\,\mathrm{d}\xi\right) = 0 .
\end{equation}

其中 $\xi$ 和 $\mathrm{d}\xi$ 都是常数，$2\pi$ 也是常数，约掉以后得到

\begin{equation}
\frac{\partial}{\partial t}\left(\rho_d R^2\right) = 0
\quad\Longrightarrow\quad
\rho_d(\xi,t)\,R(t)^2 = \rho_d(\xi,0)\,R_0^2 .
\label{eq:rhodinv}
\end{equation}

这个关系的含义是：干物质被压缩了。
半径缩小时体积变小，同样的干物质挤在更小的空间里，
所以干物质密度变大。

再进一步。初始含水率处处相同，
所以初始时刻的干物质密度 $\rho_d(\xi,0)$ 与 $\xi$ 无关，
是一个常数。代入式 \eqref{eq:rhodinv} 得到

\begin{equation}
\rho_d(\xi,t) = \rho_{d0}\frac{R_0^2}{R(t)^2} .
\label{eq:rhodconst}
\end{equation}

右边只依赖时间，不依赖 $\xi$。
这个结论非常重要：物质坐标下的干物质密度处处相同。

\why{这一点为什么关键}
因为它意味着在后面的推导中，
$\rho_d$ 可以当作常数从导数符号里提出来约掉，
方程的形式会变得非常简单。
如果 $\rho_d$ 随位置变化，方程里会多出一阶导数项，处理起来麻烦得多。

\subsection{水分守恒的完整推导}

现在对材料区域 $[0,\xi]$，也就是从圆心到这个标签之间的全部材料，
做水分衡算。

这个区域里的水分总量是多少？含水率 $C$ 的定义是每千克干物质含多少水，
所以水分总量等于

\begin{equation}
W = \int_0^\xi C\,\mathrm{d}m_d
= \int_0^\xi C\,\rho_d\,2\pi R^2\xi'\,\mathrm{d}\xi' .
\end{equation}

水分总量的变化率是

\begin{equation}
\frac{\mathrm{d}W}{\mathrm{d}t}
= \int_0^\xi \frac{\partial C}{\partial t}\,\rho_d\,2\pi R^2\xi'\,\mathrm{d}\xi' .
\label{eq:dWdt}
\end{equation}

注意这里可以把偏导数拿进积分号里面，
因为积分限 $\xi$ 是固定的材料标签，不随时间变化。
这正是用物质坐标的好处。

接下来算水分从哪里进来。

水分只能通过边界 $\xi$ 处的那个材料面进出，因为 $\xi=0$ 那端是圆心，
没有面积。

在物质坐标下，水分扩散通量怎么表示？
从半径坐标的菲克定律 $j_r = -D\dfrac{\partial C}{\partial r}$ 出发。
用链式法则换算到 $\xi$，

\begin{equation}
\frac{\partial C}{\partial r}
= \frac{\partial C}{\partial \xi}\cdot\frac{\partial \xi}{\partial r}
= \frac{\partial C}{\partial \xi}\cdot\frac{1}{R},
\end{equation}

最后一个等号用了 $\xi=r/R$，所以 $\dfrac{\partial\xi}{\partial r}=\dfrac{1}{R}$。

于是

\begin{equation}
j_r = -D\frac{\partial C}{\partial r} = -\frac{D}{R}\frac{\partial C}{\partial \xi}.
\end{equation}

注意 $j_r$ 的单位是每平方米每秒多少千克水，
它还没有考虑干物质密度的因素。
在物质坐标下做衡算时，用带干物质密度的通量更方便，也就是

\begin{equation}
j_\xi = -\frac{\rho_d D}{R}\frac{\partial C}{\partial \xi},
\end{equation}

它的含义是单位面积上、单位时间内、每单位干物质所携带的水分通量。

材料面 $\xi$ 处的面积（单位长度）是多少？
径向位置是 $r=\xi R$，周长是 $2\pi r = 2\pi\xi R$。

所以单位时间内流出这个区域的水分是

\begin{equation}
\text{流出} = j_\xi\cdot 2\pi\xi R
= -\frac{\rho_d D}{R}\frac{\partial C}{\partial \xi}\cdot 2\pi\xi R
= -2\pi\xi\rho_d D\frac{\partial C}{\partial \xi}.
\end{equation}

流入等于流出的负值，

\begin{equation}
\text{流入} = 2\pi\xi\rho_d D\frac{\partial C}{\partial \xi}.
\end{equation}

把式 \eqref{eq:dWdt} 和上式相等，

\begin{equation}
\int_0^\xi \frac{\partial C}{\partial t}\,\rho_d\,2\pi R^2\xi'\,\mathrm{d}\xi'
= 2\pi\xi\rho_d D\frac{\partial C}{\partial \xi}.
\label{eq:materialbalance}
\end{equation}

这一步就是完整的水分守恒式。

\subsection{从积分式得到微分方程}

式 \eqref{eq:materialbalance} 还是一个积分式，要化成微分方程。

做法是两边对 $\xi$ 求导。左边用微积分基本定理：
积分上限是 $\xi$，对上限求导就是被积函数在上限处的值，所以

\begin{equation}
\frac{\partial}{\partial \xi}\left[\int_0^\xi \frac{\partial C}{\partial t}\,\rho_d\,2\pi R^2\xi'\,\mathrm{d}\xi'\right]
= \frac{\partial C}{\partial t}\,\rho_d\,2\pi R^2\xi .
\end{equation}

右边求导得到

\begin{equation}
\frac{\partial}{\partial \xi}\left[2\pi\xi\rho_d D\frac{\partial C}{\partial \xi}\right]
= 2\pi\frac{\partial}{\partial \xi}\left[\xi\rho_d D\frac{\partial C}{\partial \xi}\right] .
\end{equation}

两边相等，

\begin{equation}
\frac{\partial C}{\partial t}\,\rho_d\,2\pi R^2\xi
= 2\pi\frac{\partial}{\partial \xi}\left[\xi\rho_d D\frac{\partial C}{\partial \xi}\right] .
\end{equation}

约掉两边的 $2\pi$，再把 $\rho_d$ 提出来。
这里用到式 \eqref{eq:rhodconst} 的结论，
也就是 $\rho_d$ 与 $\xi$ 无关，所以可以从导数符号里提出来，

\begin{equation}
\frac{\partial C}{\partial t}\,\rho_d\,R^2\xi
= \rho_d\frac{\partial}{\partial \xi}\left[\xi D\frac{\partial C}{\partial \xi}\right] .
\end{equation}

两边约掉 $\rho_d$，再除以 $R^2\xi$，

\begin{equation}
\boxed{\;
\left.\frac{\partial C}{\partial t}\right|_\xi
= \frac{1}{\xi R(t)^2}\frac{\partial}{\partial \xi}\left(\xi D\frac{\partial C}{\partial \xi}\right).
\;}
\label{eq:material}
\end{equation}

这就是物质坐标下的水分方程。

\subsection{换成 u 坐标}

式 \eqref{eq:material} 里还有 $\xi$，用起来和第一部分不一样。
把它换成 $u=\xi^2$，就能和前面统一。

换算关系和第五章完全一样。由 $\xi=\sqrt{u}$ 得

\begin{equation}
\frac{\partial}{\partial \xi} = 2\sqrt{u}\,\frac{\partial}{\partial u}.
\end{equation}

逐项改造式 \eqref{eq:material} 的右边。

先算括号里面，

\begin{equation}
\xi D\frac{\partial C}{\partial \xi}
= \sqrt{u}\cdot D\cdot 2\sqrt{u}\frac{\partial C}{\partial u}
= 2uD\frac{\partial C}{\partial u}.
\end{equation}

再求一次导，

\begin{equation}
\frac{\partial}{\partial \xi}\left(2uD\frac{\partial C}{\partial u}\right)
= 2\sqrt{u}\frac{\partial}{\partial u}\left(2uD\frac{\partial C}{\partial u}\right).
\end{equation}

最后除以 $\xi R^2$，

\begin{equation}
\frac{1}{\sqrt{u}R^2}\cdot 2\sqrt{u}\frac{\partial}{\partial u}\left(2uD\frac{\partial C}{\partial u}\right)
= \frac{2}{R^2}\frac{\partial}{\partial u}\left(2uD\frac{\partial C}{\partial u}\right)
= \frac{4}{R^2}\frac{\partial}{\partial u}\left(uD\frac{\partial C}{\partial u}\right).
\end{equation}

所以

\begin{equation}
\boxed{\;
\frac{\partial C}{\partial t} = \frac{4}{R(t)^2}\frac{\partial}{\partial u}\left(uD\frac{\partial C}{\partial u}\right).
\;}
\label{eq:material-u}
\end{equation}

和式 \eqref{eq:mass-u} 比较一下，两者的形式完全一样，
差别只有一个：$R$ 变成了 $R(t)$，随时间变化。

温度方程同理，

\begin{equation}
\rho c_p\frac{\partial T}{\partial t} = \frac{4}{R(t)^2}\frac{\partial}{\partial u}\left(uk\frac{\partial T}{\partial u}\right).
\end{equation}

\subsection{一个重要的结论：对流项为零}

对比一下两种写法。

在固定半径坐标下，正确的方程是要补一项的，

\begin{equation}
\left.\frac{\partial C}{\partial t}\right|_r
= \frac{1}{r}\frac{\partial}{\partial r}\left(rD\frac{\partial C}{\partial r}\right)
- \frac{\dot R}{R}\,r\,\frac{\partial C}{\partial r}.
\label{eq:radv}
\end{equation}

最后那一项叫对流项，描述干物质带着水分移动的效应。
其中 $\dot R$ 表示 $\mathrm{d}R/\mathrm{d}t$，也就是收缩速度。

在物质坐标下，这一项消失了，如式 \eqref{eq:material-u} 所示。
原因是控制体跟着材料一起走，
材料本身不迁移，所以不需要额外描述它的运动。

\why{哪一种写法更好}
物质坐标的写法更好，理由有两条。
第一，方程更短，少一个需要离散的项，
程序里不容易出错。
第二，也是更重要的，
求解区域是固定的 $0\le u\le1$，
不需要处理移动的网格。
如果用式 \eqref{eq:radv}，
每走一步网格都要重新分布，
还要处理网格移动带来的额外误差。

\why{怎么验证两种写法是等价的}
可以做一个代换检验。
设式 \eqref{eq:material-u} 成立，
把 $\xi=r/R(t)$ 代进去，
利用链式法则算出 $\left.\dfrac{\partial C}{\partial t}\right|_r$，
应该正好得到式 \eqref{eq:radv}。

具体计算如下。设 $C(r,t)=\tilde C(\xi,t)$，其中 $\xi=r/R(t)$。
固定 $r$ 对 $t$ 求导，

\begin{equation}
\left.\frac{\partial C}{\partial t}\right|_r
= \left.\frac{\partial \tilde C}{\partial t}\right|_\xi
+ \frac{\partial \tilde C}{\partial \xi}\left.\frac{\partial \xi}{\partial t}\right|_r .
\end{equation}

其中

\begin{equation}
\left.\frac{\partial \xi}{\partial t}\right|_r
= \frac{\partial}{\partial t}\left(\frac{r}{R(t)}\right)
= -\frac{r\dot R}{R^2} = -\frac{\xi\dot R}{R}.
\end{equation}

把式 \eqref{eq:material} 的右边代入第一项，
再把上面的结果代入第二项，

\begin{equation}
\left.\frac{\partial C}{\partial t}\right|_r
= \frac{1}{\xi R^2}\frac{\partial}{\partial \xi}\left(\xi D\frac{\partial \tilde C}{\partial \xi}\right)
- \frac{\xi\dot R}{R}\frac{\partial \tilde C}{\partial \xi}.
\end{equation}

现在把 $\xi$ 的导数换回 $r$ 的导数。
由 $\xi=r/R$ 得 $\dfrac{\partial}{\partial \xi}=R\dfrac{\partial}{\partial r}$。

第一项变成

\begin{equation}
\frac{1}{(r/R)R^2}\cdot R\frac{\partial}{\partial r}
\left(\frac{r}{R}D\cdot R\frac{\partial C}{\partial r}\right)
= \frac{1}{r}\frac{\partial}{\partial r}\left(rD\frac{\partial C}{\partial r}\right).
\end{equation}

第二项变成

\begin{equation}
-\frac{(r/R)\dot R}{R}\cdot R\frac{\partial C}{\partial r}
= -\frac{\dot R}{R}r\frac{\partial C}{\partial r}.
\end{equation}

两项合起来正好是式 \eqref{eq:radv}。
两种写法确实等价。

\subsection{输出表格时的坐标换算}

求解器工作在 $u$ 坐标上，而题目要求的表格列名是当前物理距离 $r$。
这两者不是一回事，必须换算。

换算关系由 $r=\xi R(t)$ 和 $u=\xi^2$ 联立得到

\begin{equation}
u = \left(\frac{r}{R(t)}\right)^2 .
\label{eq:outputmap}
\end{equation}

注意分母是当前的半径 $R(t)$，不是初始半径 $R_0$。

还有一个情况要处理：如果某一列的距离大于当前半径，
也就是 $r>R(t)$，那么这个位置已经不存在了。
例如 $r=1.5$ 厘米这一列，
当 $t>3.67$ 小时之后半径已经小于 $1.5$ 厘米，
药材外面没有东西，表格里应该留空。

\why{如果换算错了会有什么后果}
后果很严重。举一个实际例子。

在 $t=6$ 小时，半径是 $R=1.374$ 厘米。
如果要查 $r=1.0$ 厘米处的含水率。

正确的算法是

\begin{equation}
u = \left(\frac{1.0}{1.374}\right)^2 = 0.5297 .
\end{equation}

如果误用初始半径，算成

\begin{equation}
u = \left(\frac{1.0}{2.0}\right)^2 = 0.25 ,
\end{equation}

那就完全错了。$u=0.25$ 对应的是 $t=0$ 时刻位于半径 $1.0$ 厘米处的那个材料点，
而现在它已经收缩到 $1.0\times\dfrac{1.374}{2.0}=0.687$ 厘米处了。

两种算法给出的数值分别是 $1.0215$ 和 $1.3782$，
相对误差百分之三十五。

\keypoint
问题四的关键有三点。
第一，收缩时不能用固定半径坐标下的原方程，会漏掉对流项。
第二，要用物质坐标 $\xi=r/R(t)$，方程里对流项自动为零。
第三，输出表格时要用当前半径换算坐标，$r$ 超过当前半径的位置要留空。
